\section{TỔNG QUAN ĐỀ TÀI}

\subsection{Bối cảnh và động lực}
\hspace{1cm}
Trong lĩnh vực thương mại điện tử (TMĐT), dữ liệu giao dịch, hành vi duyệt web và tương
tác khách hàng được sinh ra liên tục với khối lượng lớn từ nhiều kênh bán hàng khác nhau.
Tuy nhiên, dữ liệu tác nghiệp (OLTP) thường được thiết kế để phục vụ giao dịch, không tối ưu
cho việc truy vấn phân tích đa chiều (OLAP) hay hỗ trợ ra quyết định. Điều này đặt ra nhu
cầu xây dựng một Kho dữ liệu (Data Warehouse -- DW) làm nền tảng tích hợp, làm
sạch và mô hình hóa lại dữ liệu theo hướng phân tích, từ đó phát triển các thành phần
Hệ hỗ trợ Quyết định (Decision Support System -- DSS) giúp nhà quản lý và bộ phận
kinh doanh đưa ra quyết định dựa trên dữ liệu (data-driven decision making).

\subsection{Mục tiêu}
\begin{enumerate}
    \item Phân tích yêu cầu nghiệp vụ và xác định các câu hỏi phân tích cần kho dữ liệu trả lời.
    \item Thiết kế kho dữ liệu theo mô hình đa chiều (\textit{star schema}), bao gồm quy trình
          trích xuất -- biến đổi -- nạp dữ liệu (ETL).
    \item Hiện thực một hệ thống hoàn chỉnh: tầng lưu trữ (staging + DW), tầng dịch vụ
          backend (OLAP query, API phục vụ mô hình DSS) và tầng giao diện (dashboard).
    \item Phát triển hai thành phần DSS trên nền kho dữ liệu: dự đoán rủi ro giao hàng trễ hạn
          và dự đoán rủi ro khách hàng đánh giá thấp, kèm phân tích what-if cho đơn giao trễ.
    \item Đánh giá kết quả đạt được, chỉ ra hạn chế và đề xuất hướng phát triển tiếp theo.
\end{enumerate}

\subsection{Đề tài và phạm vi}
\hspace{1cm}
Nhóm lựa chọn lĩnh vực \textbf{Thương mại điện tử} (giao thoa với \textbf{Logistics}), với bài
toán trọng tâm: xây dựng kho dữ liệu tích hợp dữ liệu đơn hàng, khách hàng, người bán, sản
phẩm, thanh toán, đánh giá và vị trí địa lý của sàn Olist (Brazil), phục vụ ba nhóm bài toán
ra quyết định:
\begin{itemize}
    \item \textbf{Phân tích doanh số và hiệu suất giao hàng} theo thời gian, ngành hàng, khu
          vực và tuyến vận chuyển (phục vụ Manager Dashboard).
    \item \textbf{Dự đoán rủi ro giao hàng trễ hạn}: chấm điểm đơn ngay khi người bán bàn giao
          cho đơn vị vận chuyển, để bộ phận vận hành can thiệp trước khi đơn bị trễ.
    \item \textbf{Dự đoán rủi ro nhận đánh giá thấp} (review $\leq$ 2 sao): chấm điểm đơn vừa
          giao tới khách, để bộ phận chăm sóc khách hàng chủ động liên hệ / tặng voucher trước
          khi khách viết đánh giá.
\end{itemize}
Hai bài toán DSS được chọn vì nhãn có sẵn trong dữ liệu thật (ngày giao thực tế so với ngày
hẹn, điểm review) và liên quan chặt chẽ với nhau: đơn giao trễ bị đánh giá $\leq$ 2 sao với tỷ
lệ 60,9\%, trong khi đơn giao đúng hẹn chỉ 9,2\%.
Phạm vi đồ án tập trung vào phần thiết kế -- hiện thực kỹ thuật (kho dữ liệu, ETL, mô hình
DSS, dashboard); dữ liệu nguồn là bộ dữ liệu mở (Kaggle), được nêu rõ ràng trong Mục~2.4;
danh sách chiến dịch chăm sóc là dữ liệu cấu hình do nhóm soạn, không tham gia phân tích. Hai mô hình DSS chỉ sử dụng dữ liệu thật.

\subsection{Đối chiếu chuẩn đầu ra học phần}
\hspace{1cm}
Bảng~\ref{tab:cdr} trình bày sự đối chiếu giữa các chuẩn đầu ra (Learning Outcomes) của học
phần CO4031 theo đề cương DCMH.CO4031.6.1 và các phần tương ứng trong báo cáo này.

\begin{table}[H]
\centering
\caption{Đối chiếu chuẩn đầu ra học phần với nội dung báo cáo}
\label{tab:cdr}
\begin{tabular}{@{}p{2cm}p{7cm}p{4.5cm}@{}}
\toprule
\textbf{CĐR} & \textbf{Mô tả} & \textbf{Mục tương ứng} \\
\midrule
L.O.1.1--1.3 & Khái niệm, vai trò, ví dụ thực tế về DW và DSS & Mục 1, Mục 5.1 \\
L.O.2.1--2.3 & Quy trình, nguyên lý thiết kế và hiện thực DW & Mục 3, Mục 4 \\
L.O.3.1--3.3 & Khái niệm MSS/DSS, các giai đoạn ra quyết định, mô hình hóa & Mục 5.1 \\
L.O.4.1--4.2 & Nguyên tắc phân tích, thiết kế và hiện thực DSS đơn giản & Mục 5 \\
\bottomrule
\end{tabular}
\end{table}
